% !TITLE 迢迢牵牛星
% !AUTHOR 无名氏
% !DYNASTY 两汉
% !GENRE 五言古诗
% !ORDER 7
% !SOURCE 《古诗十九首》；收于《文选》
% !EDITION 《古诗十九首集释》，隋树森，中华书局（2020年版为候选）
% !EDITION_STATUS pending

\begin{poem}
迢迢牵牛星\textsuperscript{1}，皎皎河汉女\textsuperscript{2}。
纤纤擢素手\textsuperscript{3}，札札弄机杼\textsuperscript{4}。
终日不成章\textsuperscript{5}，泣涕零如雨\textsuperscript{6}。
河汉清且浅，相去复几许\textsuperscript{7}？
盈盈一水间\textsuperscript{8}，脉脉不得语\textsuperscript{9}。
\end{poem}

\begin{annotations}
\notetext{1}{迢迢：遥远的样子。}
\notetext{2}{皎皎：明亮洁白。河汉女：指织女星。河汉，银河。}
\notetext{3}{纤纤：细柔的样子。擢：伸出、举起。素手：洁白的手。}
\notetext{4}{札札：织机的声音。机杼：织布机。杼，织布的梭子。}
\notetext{5}{章：布幅上的纹理，这里指织成整幅的布。}
\notetext{6}{零：落下。}
\notetext{7}{几许：多少，这里指多远。}
\notetext{8}{盈盈：水清浅充盈的样子。一水：指银河。}
\notetext{9}{脉脉：含情相视的样子。}
\end{annotations}

\begin{appreciation}
《迢迢牵牛星》是古诗十九首的第十首，也是这组诗里写相思写得最含蓄的之一。

诗一开口就把镜头推向星空：“迢迢牵牛星，皎皎河汉女。”牵牛星那么遥远，银河边的织女那么明亮。两个叠字，一远一亮，先把天上的一对恋人摆在我们眼前。可接下来，镜头忽然落到了人间的一架织机前：“纤纤擢素手，札札弄机杼。”一双纤细白净的手伸出来，织机札札地响。

这两句写得极细。纤纤是手的形状，素是手的颜色，擢是手抬起来的动作，札札是机杼的声音。读到这里，你几乎能看见那个女子坐在织机前，也能听见那台织机。可是“终日不成章”——织了一整天，一匹布也没织成。为什么？“泣涕零如雨”，眼泪像雨一样落下来。

织女的工作是织云锦，天上那一片片云霞都是她织的。这样一双巧手，此刻却连一匹布也织不下去。诗里没有写她在想谁，只写她的手在动、机在响、泪在落，相思的重量全压在这几个动作上。

“河汉清且浅，相去复几许？”银河又清又浅，两岸相隔能有多远呢？这一问看似轻巧，其实最狠。如果隔着千山万水，倒也罢了；偏偏只隔一条清浅的水，抬脚就能过去，可就是过不去。“盈盈一水间，脉脉不得语”——水光盈盈，只此一水之隔，两个人含情相望，一句话也说不出来。

世上的阻隔大多如此。真正让人难受的，往往不是天涯海角，而是近在眼前却够不着。一道浅浅的银河，把“相去复几许”变成了永远。

牵牛织女作为星名，最早见于《诗经·大东》：“跂彼织女，终日七襄”“睆彼牵牛，不以服箱”，那时它们还只是两颗名不副实的星星。到了这首诗，星星有了人形，有了手，有了眼泪，有了说不出口的话。再往后，就有了七夕、鹊桥、乞巧，有了秦观那句“两情若是久长时，又岂在朝朝暮暮”。

从两颗星到一对夫妻，中间隔着许多代人的想象，这首诗是其中极要紧的一步。

天上的银河还在。人间的银河，换了一条又一条。
\end{appreciation}
